\input{papers/template/style.tex}

\title{Research Directions: Time, Complexity, Causality, and Learning
Theory}
\date{}

\begin{document}
  \maketitle

  \section*{Promising Research Directions in Machine Learning and AI}

  Drawing on my 15 years of experience across finance, predictive
  maintenance, energy forecasting, and complex real-world systems, I believe
  the next frontier of machine learning might lie not in scaling existing
  paradigms, but in rethinking their foundations.

  \begin{itemize}
    \item \textbf{Reconceptualizing Time in Learning Systems}

      Modern learning theory largely treats time as an external index over
      data. Yet in real systems, time is not merely a coordinate—it is generative,
      causal, and structural. Markets evolve, machines degrade, energy
      systems adapt, and agents respond to predictions.

      A fundamental question emerges: \emph{Is time truly special, or have
      we artificially separated it from the feature space due to historical
      modeling choices?}

      Is it possible a unification in which time is embedded directly into
      the representation of the system—where features, dynamics, and
      evolution are treated within a single formal structure? Such a framework
      would remove the boundary between static learning and dynamical systems,
      leading to new principles for:
      \begin{itemize}
        \item Learning in non-stationary and path-dependent environments,

        \item Modeling feedback loops between predictions and reality,

        \item Understanding adaptation as an inherent property of learning
          systems.
      \end{itemize}

      This shift would move machine learning from pattern recognition
      toward \emph{temporal reasoning about evolving worlds}.

    \item \textbf{A Theory of Complexity for Time, Causality, and
      Discovery}

      Classical notions of complexity—such as Minimum Description Length
      (MDL) and VC dimension—were developed for static hypothesis classes
      under i.i.d.\ assumptions. However, real-world learning unfolds in
      time, under causal constraints, and often in interaction with the
      environment.

      We need a new theory of complexity that captures:
      \begin{itemize}
        \item Capacity in time-evolving and feedback-driven systems,

        \item The cost of modeling causal structure rather than mere correlation,

        \item The adaptive nature of the research and model selection process
          itself.
      \end{itemize}

      Such a theory would quantify not only the complexity of models, but
      the complexity of \emph{learning trajectories}. It would connect generalization,
      stability, causality, and scientific discovery within a unified
      mathematical framework.

      Ultimately, this line of research aims to answer a deeper question:
      \emph{What does it mean to learn in a world that changes in
      response to being modeled?}
  \end{itemize}

% ##############################################################################
  \section{Time as a Feature of Machine Learning}

  \subsection*{Core Idea}

  Treat time not only as an index but as a structural feature affecting hypothesis
  complexity and generalization.

  Define a time-indexed hypothesis class:
  \[
    \mathcal{H}_{t}= \{h_{t}(x) : t \in \mathbb{T}\}
  \]

  Assume bounded drift:
  \[
    \|h_{t+1}- h_{t}\| \le \epsilon
  \]

  Generalize with drift

  Classical PAC bound:
  \[
    n \gtrsim \frac{VC(\mathcal{H}) + \log(1/\delta)}{\epsilon^{2}}
  \]

  Time-aware extension:
  \[
    n \gtrsim \frac{VC(\mathcal{H}) + D_{T}+ \log(1/\delta)}{\epsilon^{2}}
  \]

  where $D_{T}$ measures cumulative drift:
  \[
    D_{T}= \sum_{t=1}^{T-1}\|h_{t+1}- h_{t}\|
  \]

  \textbf{Examples}

  \begin{itemize}
    \item Stock market prediction: $h_{t}(x)$ predicts returns from
      features where the relationship changes due to market regimes (bull/bear
      markets, crisis periods)

    \item Adaptive recommendation systems: User preferences drift over
      time, requiring hypothesis updates $h_{t}\to h_{t+1}$ as tastes evolve

    \item Climate modeling: Physical relationships may shift due to
      anthropogenic forcing, violating stationarity assumptions

    \item Medical diagnosis: Disease patterns and diagnostic criteria
      evolve with new treatments and emerging pathogens

    \item Natural language processing: Language models trained on pre-2020
      data fail to understand pandemic-related terminology or evolving
      slang. The relationship between text features and meaning changes as
      vocabulary and usage patterns drift.

    \item Fraud detection: Fraudsters continuously adapt their
      strategies to evade detection, creating an adversarial drift where
      $P(y|x, t)$ is actively manipulated by intelligent agents.

    \item Energy demand forecasting: The relationship between
      temperature and electricity usage changes as solar panel adoption
      increases and electric vehicles become common, fundamentally altering
      load patterns.
  \end{itemize}

  \textbf{Provocative Questions}

  \begin{itemize}
    \item Can a model that constantly adapts to time still learn anything
      generalizable? If $h_{t}$ changes at every timestep, is this learning
      or mere tracking?

    \item Is forgetting necessary for generalization in non-stationary environments?
      Classical theory rewards more data, but in drifting environments,
      old data may hurt performance.

    \item Does the notion of "ground truth" even make sense in time-varying
      systems? If $y = f_{t}(x)$ where $f_{t}$ itself evolves, what are we
      actually trying to learn—the current $f_{t}$ or the meta-function that
      generates the sequence $\{f_{t}\}$?

    \item Can we have PAC-style guarantees without stationarity? Or does
      non-stationarity fundamentally break the connection between
      empirical and true risk?

    \item Should we penalize model complexity or model \textit{stability}?
      A complex but stable model might generalize better over time than a
      simple but volatile one.

    \item Is the distinction between "learning from data" and "tracking
      a signal" just a matter of timescale? At what rate of change does
      learning become impossible?
  \end{itemize}

  \textbf{Research Topics}

  \begin{itemize}
    \item VC dimension of time-indexed hypothesis classes

    \item Time-dependent Rademacher complexity

    \item Stability under evolving distributions

    \item Meta-learning for non-stationary tasks
  \end{itemize}

% ##############################################################################
  \section{Quasi-Stationary Learning}

  \subsection*{Types of Non-Stationarity}

  There are several types of non-stationarities

  \begin{itemize}
    \item \textbf{Covariate shift:} $P(x)$ changes
      \begin{itemize}
        \item Example: Training spam filter on 2020 emails, deploying on
          2025 emails (vocabulary changes, but spam concept remains)
      \end{itemize}

    \item \textbf{Concept drift:} $P(y|x)$ changes
      \begin{itemize}
        \item Example: Credit scoring where same financial profile means
          different risk due to changing economic conditions
      \end{itemize}

    \item \textbf{Label shift:} $P(y)$ changes
      \begin{itemize}
        \item Example: Disease prevalence changes during epidemic (features
          stay same, base rates change)
      \end{itemize}

    \item \textbf{Regime shift:} abrupt structural change
      \begin{itemize}
        \item Example: COVID-19 causing discontinuous change in consumer
          behavior, market dynamics, work patterns
      \end{itemize}
  \end{itemize}

  \subsection*{Temporal VC Dimension}

  Define:
  \[
    VC_{T}= VC\left(\bigcup_{t=1}^{T}\mathcal{H}_{t}\right)
  \]

  \textbf{Examples:}
  \begin{itemize}
    \item Consider a linear classifier $h_{t}(x) = \text{sign}(w_{t}^{\top}
      x)$ where $w_{t}$ drifts smoothly. At any fixed $t$, $VC(\mathcal{H}
      ) = d+1$ (standard for linear classifiers in $\mathbb{R}^{d}$).
      But over $T$ timesteps with unbounded drift, $VC_{T}$ can grow without
      bound as the union captures increasingly complex decision boundaries.

    \item A neural network with fixed architecture but time-varying weights.
      If weights can change arbitrarily, $VC_{T}$ may equal the VC dimension
      of the universal function approximator class, even if each $h_{t}$
      individually has bounded complexity.

    \item Seasonal models: A retailer uses $h_{\text{holiday}}$ during
      December and $h_{\text{regular}}$ otherwise. Then $VC_{T}$ is at least
      $\max(VC(h_{\text{holiday}}), VC(h_{\text{regular}}))$ but could be
      larger if the union creates new decision boundaries.
  \end{itemize}

  Questions:
  \begin{itemize}
    \item Is complexity additive over time?

    \item Can we define a temporal VC dimension?

    \item If $VC_{T}\to \infty$ as $T \to \infty$, does this mean all
      long-term predictions are impossible?

    \item Can two models have identical $VC(\mathcal{H})$ but vastly different
      $VC_{T}$ due to different drift patterns?

    \item Does smooth drift (small $\|h_{t+1}- h_{t}\|$) lead to sublinear
      growth in $VC_{T}$, or is any non-zero drift enough to cause
      unbounded complexity?

    \item If we observe concept drift but don't model it, does the
      effective $VC_{T}$ perceived by a stationary learner explode?
  \end{itemize}

  \subsection*{Kolmogorov Complexity Over Time}

  Measure compressibility of a data stream:
  \[
    K(x_{1:T})
  \]

  Incremental description length:
  \[
    L = \sum_{t=1}^{T}K(x_{t}\mid x_{<t})
  \]

  Use compression rate changes to detect regime shifts.

  \textbf{Example:} A financial time series that suddenly becomes less
  compressible (higher $K(x_{t}\mid x_{<t})$) may indicate a regime
  change where past patterns no longer predict future behavior—think of
  the 2008 crisis or COVID-19 market turbulence.

  \textbf{Provocative Question:} \textit{Is there a fundamental limit to
  predictability based on Kolmogorov complexity? If $K(x_{t}\mid x_{<t})$
  approaches $|x_{t}|$ (data is incompressible given past), does this
  imply the system has become random or that our model class is too weak?}

  \subsection*{Forgetting Mechanisms}

  Exponential decay weights:
  \[
    w_{t}= e^{-t/\tau}
  \]

  \textbf{Examples:}
  \begin{itemize}
    \item Reinforcement learning in changing environments: An RL agent
      playing a game where rules gradually change must weight recent
      experience more heavily. Too much memory and it keeps using obsolete
      strategies; too little and it can't learn patterns.

    \item Online advertising: Click-through rate models must balance
      learning user preferences (requires memory) with adapting to
      changing tastes (requires forgetting). A user who clicked sports ads
      in January might prefer fashion ads by March.

    \item Robotics in aging systems: A robot learning to compensate for
      motor degradation must forget its old calibration while retaining
      general manipulation skills. Which memories should persist and which
      should decay?
  \end{itemize}

  Research:
  \begin{itemize}
    \item Optimal $\tau$ under bounded drift

    \item Bias-variance tradeoff in non-stationary environments

    \item Regret bounds with exponential forgetting
  \end{itemize}

  \textbf{Questions:}
  \begin{itemize}
    \item If we forget too fast ($\tau$ too small), we lose
      generalization. If we forget too slow, we overfit to obsolete data.
      Is there a "no free lunch" theorem for memory decay?

    \item Should forgetting rate $\tau$ itself be learned from data? Or
      does this create a "meta-overfitting" problem where we overfit to
      recent drift patterns?

    \item Can we design forgetting mechanisms that are "content-aware"—
      forgetting irrelevant noise while retaining stable patterns? Is this
      even possible without knowing the future?

    \item Does human memory's selective forgetting (we remember
      important events but forget mundane details) suggest an optimal forgetting
      strategy for ML? Can we formalize "importance-weighted forgetting"?

    \item Is there a fundamental trade-off between adaptation speed and stability?
      Can we prove lower bounds on regret in terms of drift rate and
      memory window?
  \end{itemize}

  \subsection*{VC Dimension of Causal / Bayesian Networks}

  Complexity depends on:
  \begin{itemize}
    \item Number of nodes

    \item Maximum indegree

    \item Parametric family of conditionals
  \end{itemize}

  Research:
  \begin{itemize}
    \item Structural VC dimension of DAGs

    \item Sample complexity of causal discovery

    \item MDL penalties for structure learning
  \end{itemize}

  \textbf{Example:} A 10-node causal DAG with max indegree 3 and linear
  Gaussian conditionals has far lower VC dimension than a fully-connected
  10-node network. This suggests causal structure provides inductive bias
  that improves generalization—but only if the assumed structure is correct.

  \textbf{Provocative Questions:}
  \begin{itemize}
    \item Does learning the wrong causal structure hurt generalization more
      than ignoring causality entirely and using pure correlation?

    \item Can we define a "causal VC dimension" that captures not just the
      complexity of the model but the complexity of interventions it can
      represent?

    \item If two causal graphs are Markov equivalent (indistinguishable
      from observational data), do they have the same VC dimension?

    \item Is there a fundamental trade-off between causal interpretability
      and predictive accuracy?
  \end{itemize}

% ##############################################################################
  \section{MDL Extensions}

  \subsection*{Including the Research Process}

  Total complexity:
  \[
    C = C_{\text{model}}+ C_{\text{search}}
  \]

  Study:
  \begin{itemize}
    \item VC dimension of hyperparameter search

    \item Complexity of AutoML pipelines

    \item Generalization with architecture search
  \end{itemize}

  \textbf{Example:} Consider AutoML trying 1000 different model
  architectures, each with 100 hyperparameter combinations. Even if the
  "best" model has low VC dimension, the search process itself adds $\log
  (1000 \times 100) \approx 17$ bits of complexity. This hidden cost
  explains why AutoML often overfits despite using "simple" final models.

  \textbf{Provocative Questions:}
  \begin{itemize}
    \item If a researcher tries 100 approaches but only publishes the best
      one, the community systematically underestimating model complexity?
      How to fix this?

    \item Can we detect overfitting from the research process itself by analyzing
      the "trajectory" through model space rather than just the final model?

    \item Should scientific venues require a "complexity tax" where papers
      must report $C_{\text{search}}$ alongside model accuracy? Would
      this reveal that many "breakthroughs" are actually just
      overfitting?

    \item If two researchers independently discover the same solution,
      does this reduce its effective complexity? Does convergent discovery
      serve as evidence against overfitting?

    \item Can we formalize the idea that "surprising" results (low prior
      probability) should be penalized more heavily for multiple testing
      than "expected" results?

    \item Is transfer learning a way to "amortize" search cost across tasks?
      If a pre-trained model is fine-tuned with minimal search, does this
      reduce $C_{\text{search}}$ for the downstream task?
  \end{itemize}

  \subsection*{Backtesting Complexity}

  Effective complexity:
  \[
    VC_{\text{eff}}= VC(\mathcal{H}) + \log(N_{\text{strategies tested}})
  \]

  Research:
  \begin{itemize}
    \item VC dimension of trading strategies

    \item Multiple testing corrections via learning theory

    \item Overfitting detection using Rademacher complexity
  \end{itemize}

  \textbf{Example:} A hedge fund backtests 10,000 trading strategies on
  20 years of data and selects the top performer. Even if that strategy has
  Sharpe ratio 3.0 in-sample, the effective VC dimension includes the
  search over 10,000 strategies, dramatically inflating the risk of overfitting.
  The "published" strategy may have $VC_{\text{eff}}\gg VC(\mathcal{H})$.

  \textbf{Real-world scenario:} Quantitative trading firms often find that
  strategies with impressive backtest performance degrade immediately upon
  deployment—a phenomenon directly explained by ignoring $C_{\text{search}}$
  in complexity estimates.

  \textbf{Provocative Questions:}
  \begin{itemize}
    \item Is the "replication crisis" in science fundamentally a failure
      to account for $C_{\text{search}}$ in reported p-values?

    \item Can we use learning theory to derive optimal publication policies?
      For instance, should journals require authors to report how many analyses
      they attempted?

    \item If a strategy was discovered by "accident" (zero search cost) versus
      intensive backtesting (high search cost), should we trust the
      accidental discovery more?

    \item Does cross-validation protect against overfitting the research
      process itself, or does it only protect against overfitting within
      a single model?

    \item If we condition on "this backtest passed," we've selected from
      a larger hypothesis class than we realize—can we formalize this
      selection bias using VC theory?
  \end{itemize}
\end{document}
